\section{Phase 1: Foundations and Controlled Medical Testbed}

Phase~1 establishes the foundation for evaluating whether compositional generative
modelling can behave in a \textit{semantically meaningful} way in complex real-world
settings. This phase directly advances \textbf{Research Goal~1} and provides the
first empirical test of \textbf{Hypothesis~H1}, which states that composition is
semantically meaningful only when the underlying models share a sufficiently aligned
and compatible representation of the generative factors they encode.

A key contribution of this phase is the introduction of a clinically realistic
evaluation environment based on TB and Silicosis. These diseases provide an ideal
testbed for probing semantic compositionality because they exhibit overlapping
patterns, affect the same anatomical regions, and interact in ways that violate
common simplifying assumptions used in many compositional methods. Instead of relying
on toy datasets with artificially disentangled factors, we use a domain where the
true generative structure is complex, entangled, and clinically meaningful. This
setting enables us to reveal when different compositional operators genuinely reflect
semantic combination and when they fail, despite being mathematically well-defined.

\paragraph{Score-Based Composition.}
We evaluate score-based operators by training TB and Silicosis diffusion models in a
shared latent space and combining their scores linearly. This tests whether adding
gradient fields can produce coherent mixed pathology. Success would indicate that
latent alignment and mild factor compatibility are enough for meaningful hybrid
images. Failure—manifesting as blended textures, misplaced lesions, or anatomically
inconsistent shapes—indicates that score addition does not correspond to semantic
composition when the underlying factors overlap or are entangled.

\paragraph{Energy-Based PoE Composition.}
We assess Product-of-Experts composition by summing the log-densities of the TB and
Silicosis models. In principle, this corresponds to a logical “AND.” In practice,
PoE is sensitive to conflicts between models. Because TB and Silicosis share
anatomical regions and exhibit similar textural signatures, PoE may either suppress
valid modes or amplify incompatible ones. Observing whether PoE yields coherent
co-morbidity or collapses into unrealistic hybrids helps reveal when independence
assumptions break down and how this affects semantic compositionality.

\paragraph{Density-Based Composition (SuperDiff).}
We apply SuperDiff to compute per-model log-densities along a shared trajectory and
form composite images using equal-density or mixture rules. Density-based operators
offer the most principled mathematical formulation; however, mathematical correctness
does not guarantee semantic correctness. Because the log-density lives in pixel
space—not in pathology factor space—equal-density constraints may or may not
correspond to meaningful co-morbidity. Radiological interpretation determines whether
the composed density reflects realistic combinations of TB and Silicosis or produces
artifacts arising from misaligned latent structure.

\paragraph{Cross-Operator Comparative Evaluation.}
By applying all three classes of compositional methods to the same TB and Silicosis
models, we obtain a unified view of their behaviour under identical structural
conditions. This allows us to identify when composition succeeds in capturing real
co-morbidity and when it fails due to factor entanglement, latent misalignment, or
geometric incompatibility. Success cases highlight the structural conditions under
which semantic composition is achievable. Failure cases provide insight into the
limits of current methods and motivate the refinements explored in later phases.

Phase~1 therefore provides a direct, clinically meaningful assessment of
compositional generative modelling. By grounding our evaluation in a realistic and
structurally challenging domain, we identify the conditions under which compositional
operators produce plausible and interpretable results and the conditions under which
they break down. These findings establish the necessary foundations for the
subsequent phases of the research.
